\section{De la recta al plano y al espacio}

En los capítulos anteriores describimos el movimiento mediante una coordenada:
la posición sobre una carretera recta o la altura de un cuerpo en movimiento
vertical. Las definiciones de desplazamiento, velocidad y aceleración siguen
siendo válidas cuando la trayectoria deja de ser recta. Lo que cambia es la
información necesaria para indicar la posición y la orientación del movimiento.

Pensemos en una cámara transportada por un vehículo aéreo. Para localizarla
no basta conocer su altura: también necesitamos dos coordenadas horizontales.
Si mantiene su altura, su movimiento puede estudiarse en un plano. Si asciende
mientras cambia de posición horizontal, necesitamos en general una descripción
espacial. Podemos representarla como una partícula si interesa su ubicación
y no su tamaño ni su orientación.

La pregunta central es cómo combinar los cambios de varias coordenadas para
describir \emph{un mismo movimiento}. Utilizaremos vectores y recuperaremos,
sobre cada eje, las relaciones cinemáticas conocidas. Este capítulo establece
el lenguaje general; proyectiles, movimiento circular y comparación entre
observadores se desarrollarán en los capítulos siguientes.

\section{Referencia, coordenadas y vector de posición}

\subsection{Una referencia fija para toda la descripción}

Elegimos un origen espacial \(O\), tres ejes cartesianos perpendiculares y un
reloj con un origen temporal definido. Mantendremos los ejes fijos respecto
del observador elegido, sin cambiar su orientación durante el análisis.
Los unitarios \(\ihat\), \(\jhat\) y \(\khat\) apuntan en los sentidos
positivos de \(x\), \(y\) y \(z\), respectivamente, y son adimensionales.
Usaremos un sistema derecho, con \(\ihat\pcross\jhat=\khat\).

Cuando el plano representado sea vertical, tomaremos habitualmente \(+x\)
hacia la derecha y \(+y\) hacia arriba, como en el capítulo anterior;
\(+z\) saldrá del plano de la página. Para un plano horizontal indicaremos
las direcciones geográficas o de laboratorio correspondientes. Ninguna letra
representa por sí misma una dirección física universal.

Las coordenadas se expresan en metros y el tiempo en segundos. Denotamos
el instante inicial mediante \(t_0\); el subíndice \(0\) indicará valores
en ese instante. El origen espacial no tiene que coincidir con la posición inicial.

\subsection{El vector y sus componentes}

\begin{definicion}[Vector de posición]
El vector de posición \(\vect{r}(t)\) va desde el origen hasta la posición
de la partícula en el instante \(t\). En coordenadas cartesianas,
\[ \vect{r}(t)=x(t)\ihat+y(t)\jhat+z(t)\khat. \]
Las coordenadas \(x(t)\), \(y(t)\) y \(z(t)\) son sus componentes escalares.
\end{definicion}

Distinguimos \(x\), componente escalar con signo; \(x\ihat\), contribución
vectorial sobre el eje \(x\); y \(\vect{r}\), vector completo. Su magnitud,
\[ r=\norm{\vect{r}}=\sqrt{x^2+y^2+z^2}, \]
es la distancia al origen, no la distancia recorrida. Una coordenada negativa
indica una ubicación respecto del origen y no un sentido de movimiento.

Para un movimiento contenido en el plano \(xy\), elegimos el origen en ese
plano, de modo que \(z(t)=0\) y
\[ \vect{r}(t)=x(t)\ihat+y(t)\jhat. \]
Si el plano es paralelo a \(xy\), pero no contiene el origen, \(z\) será
constante y no necesariamente nula. Un movimiento rectilíneo oblicuo puede
tener tres coordenadas variables: que varíen tres coordenadas no demuestra
que la trayectoria sea una curva espacial. Podemos reducir la descripción
si elegimos ejes adecuados a la geometría.

\begin{figure}[htbp]
    \centering
    \begin{tikzpicture}[x={(0.95cm,-0.32cm)},y={(0cm,1cm)},z={(-0.7cm,-0.5cm)}]
        \coordinate (O) at (0,0,0);
        \coordinate (P) at (3,2.5,1.5);
        \draw[-{Latex},thick] (O) -- (4.4,0,0) node[right] {$x$};
        \draw[-{Latex},thick] (O) -- (0,3.8,0) node[above] {$y$};
        \draw[-{Latex},thick] (O) -- (0,0,3) node[left] {$z$};
        \draw[dashed,GrisPrincipal] (3,0,0) -- (3,0,1.5) -- (0,0,1.5);
        \draw[dashed,GrisPrincipal] (3,0,1.5) -- (P) -- (0,2.5,1.5) -- (0,2.5,0);
        \draw[-{Latex},very thick,AzulPrincipal] (O) -- (P)
            node[midway,right] {$\vect{r}$};
        \fill (P) circle (2pt) node[above right] {$P(x,y,z)$};
        \node[above left] at (O) {$O$};
        \node[below] at (3,0,0) {$x$};
        \node[left] at (0,2.5,0) {$y$};
        \node[below left] at (0,0,1.5) {$z$};
    \end{tikzpicture}
    \caption{Sistema cartesiano en perspectiva. Las líneas auxiliares
    localizan las coordenadas de \(P\); los tres ejes físicos son perpendiculares.}
\end{figure}

\subsection{Trayectoria y ecuaciones paramétricas}

Las funciones \(x(t)\), \(y(t)\) y \(z(t)\) son las \textbf{ecuaciones
paramétricas} del movimiento: el mismo valor de \(t\) selecciona las tres
coordenadas de una única posición. Al variar el tiempo, el extremo de
\(\vect{r}(t)\) dibuja la trayectoria.

A veces podemos eliminar \(t\) para obtener una relación entre coordenadas,
como \(y=f(x)\). Esta describe el camino geométrico, pero por sí sola no
dice cuándo pasa la partícula por cada punto, con qué rapidez ni en qué
sentido. Tampoco toda trayectoria admite una única función \(y=f(x)\):
una misma coordenada \(x\) puede corresponder a varias posiciones.

\begin{ejemplo}[La trayectoria no determina el ritmo del movimiento]
En una mesa horizontal fijamos \(+x\) hacia la derecha y \(+y\) hacia el
fondo. Dos puntos móviles parten del origen en \(t=0\), con
\[
\begin{aligned}
 x_A(t)&=(\SI{1}{\meter\per\second})t,
 &y_A(t)&=(\SI{2}{\meter\per\second})t,\\
 x_B(t)&=(\SI{1}{\meter\per\second\squared})t^2,
 &y_B(t)&=(\SI{2}{\meter\per\second\squared})t^2.
\end{aligned}
\]
Para \(t\geq0\), ambos recorren la semirrecta \(y=2x\), \(x\geq0\).
A recorre incrementos iguales de posición en tiempos iguales; B recorre
incrementos cada vez mayores. Coinciden en
\((\SI{1}{\meter},\SI{2}{\meter})\) a \(t=\SI{1}{\second}\), pero no
tienen el mismo movimiento. Cada término tiene unidades de longitud:
los coeficientes de \(t\) y \(t^2\) deben tener unidades distintas.
\end{ejemplo}

\section{Desplazamiento y distancia en una trayectoria curva}

Entre \(t_i\) y \(t_f>t_i\), el cambio de posición es
\[ \Delta\vect{r}=\vect{r}_f-\vect{r}_i
=\Delta x\ihat+\Delta y\jhat+\Delta z\khat, \]
donde \(\Delta x=x_f-x_i\), \(\Delta y=y_f-y_i\) y
\(\Delta z=z_f-z_i\). Restamos componentes del mismo eje.
Geométricamente, \(\Delta\vect{r}\) va desde la posición inicial hasta la
final; no comienza necesariamente en el origen ni sigue la curva recorrida.

Su magnitud es
\[ \norm{\Delta\vect{r}}=\sqrt{(\Delta x)^2+(\Delta y)^2+(\Delta z)^2}. \]
La distancia recorrida \(s\), introducida en el capítulo 7, suma las
longitudes de todos los tramos efectivamente recorridos. Por tanto,
\[ s\geq\norm{\Delta\vect{r}}. \]
Para un desplazamiento no nulo, hay igualdad si se recorre el segmento recto
entre los extremos sin invertir el sentido; las pausas no alteran esa
igualdad. Un camino curvo puede producir desigualdad estricta aunque la
partícula nunca retroceda sobre su trayectoria.

\begin{figure}[htbp]
    \centering
    \begin{tikzpicture}[scale=0.95]
        \draw[-{Latex}] (0,0) -- (6.7,0) node[right] {$x$};
        \draw[-{Latex}] (0,0) -- (0,4.3) node[above] {$y$};
        \node[below left] at (0,0) {$O$};
        \coordinate (A) at (1.3,1);
        \coordinate (B) at (5.5,2.8);
        \draw[thick,GrisPrincipal] (A) .. controls (1.7,3.8) and (4.2,4.1) .. (B);
        \node[above] at (3.3,3.4) {trayectoria};
        \draw[-{Latex},thick,AzulPrincipal] (0,0) -- (A)
            node[midway,below right] {$\vect{r}_i$};
        \draw[-{Latex},thick,AzulPrincipal] (0,0) -- (B)
            node[pos=0.64,below right] {$\vect{r}_f$};
        \draw[-{Latex},very thick,VerdePrincipal] (A) -- (B)
            node[midway,above] {$\Delta\vect{r}$};
        \fill (A) circle (2pt) node[left] {$P_i$};
        \fill (B) circle (2pt) node[right] {$P_f$};
    \end{tikzpicture}
    \caption{El desplazamiento une los extremos del recorrido.
    La suma \(\vect{r}_i+\Delta\vect{r}=\vect{r}_f\) no depende
    de la forma del camino. Esquema sin escala.}
\end{figure}

\begin{ejemplo}[Dos tramos perpendiculares]
Un vehículo de inspección se mueve sobre un suelo horizontal. Tomamos
el origen en una esquina, \(+x\) al este y \(+y\) al norte.
Parte de \((\SI{1}{\meter},\SI{2}{\meter})\), avanza \(\SI{3}{\meter}\)
al este y después \(\SI{4}{\meter}\) al norte. Su posición final es
\((\SI{4}{\meter},\SI{6}{\meter})\), de modo que
\[ \Delta\vect{r}=(\SI{3}{\meter})\ihat+(\SI{4}{\meter})\jhat,
\qquad \norm{\Delta\vect{r}}=\SI{5}{\meter}. \]
La distancia recorrida es \(s=\SI{7}{\meter}\). La distancia final al origen
es otra cantidad: \(r_f=\sqrt{4^2+6^2}\,\si{\meter}\approx\SI{7.21}{\meter}\).
Los tres resultados responden preguntas distintas. El recorrido por tramos
idealiza el giro; no describe la aceleración durante una esquina instantánea.
\end{ejemplo}

\section{Velocidad vectorial y rapidez}

\subsection{Velocidad media}

Generalizamos la definición del capítulo 8 sin cambiar su significado:
\[ \overline{\vect{v}}=\frac{\Delta\vect{r}}{\Delta t},
\qquad \Delta t=t_f-t_i>0. \]
Al dividir por un escalar positivo, la velocidad media tiene la dirección
y el sentido del desplazamiento cuando este no es nulo. Sus componentes son
\[ \bar v_x=\frac{\Delta x}{\Delta t},\qquad
\bar v_y=\frac{\Delta y}{\Delta t},\qquad
\bar v_z=\frac{\Delta z}{\Delta t}. \]
El vector y cada componente tienen unidades de \(\si{\meter\per\second}\).
La velocidad media no tiene por qué ser tangente al camino en sus extremos.
Si se vuelve al punto inicial, es nula aunque haya habido movimiento.

La rapidez media conserva la definición \(s/\Delta t\). Por ello,
\[ \text{rapidez media}=\frac{s}{\Delta t}
\geq\norm{\overline{\vect{v}}}. \]
Si el recorrido anterior dura \(\SI{5}{\second}\),
\(\overline{\vect{v}}=(\SI{0.6}{\meter\per\second})\ihat
+(\SI{0.8}{\meter\per\second})\jhat\). Su magnitud es
\(\SI{1}{\meter\per\second}\), mientras la rapidez media es
\(\SI{1.4}{\meter\per\second}\). Estos datos no determinan las velocidades
instantáneas durante cada tramo.

\subsection{Velocidad instantánea y tangente a la trayectoria}

Al reducir el intervalo alrededor de un instante, la secante que une las
posiciones sucesivas se aproxima a la tangente de una trayectoria suave.
Así se obtiene la generalización vectorial de la derivada ya conocida:
\begin{definicion}[Velocidad instantánea]
\[ \vect{v}(t)=\lim_{\Delta t\to0}
\frac{\vect{r}(t+\Delta t)-\vect{r}(t)}{\Delta t}
=\deriv{\vect{r}}{t}. \]
Cuando \(\vect{v}\ne\vect{0}\), su dirección es tangente a la trayectoria
y su sentido es el del movimiento.
\end{definicion}

Como los unitarios de nuestros ejes fijos no cambian con el tiempo,
derivar el vector equivale a derivar sus componentes:
\[ \vect{v}=v_x\ihat+v_y\jhat+v_z\khat,
\qquad v_x=\deriv{x}{t},\quad v_y=\deriv{y}{t},\quad v_z=\deriv{z}{t}. \]
Cada componente expresa la tasa de cambio de una coordenada. Si \(v_x<0\),
\(x\) disminuye en ese instante, aunque \(x\) sea positivo. Si \(v_x=0\),
la proyección sobre \(x\) está momentáneamente en reposo, pero la partícula
puede seguir moviéndose sobre los otros ejes. El reposo instantáneo exige
que todas las componentes sean simultáneamente nulas.

La rapidez es la magnitud de la velocidad:
\[ v=\norm{\vect{v}}=\sqrt{v_x^2+v_y^2+v_z^2}\geq0. \]
No se obtiene sumando componentes ni sus valores absolutos. Tampoco es,
en general, \(\dd r/\dd t\): una partícula puede cambiar de posición sin
variar su distancia al origen. Cuando \(\vect{v}=\vect{0}\), el vector
velocidad no tiene dirección definida; la frase ``velocidad tangente'' no
permite asignarle una orientación a un vector nulo.

\subsection{Orientación de la velocidad}

En el plano \(xy\), para \(v>0\), medimos \(\theta_v\) desde \(+x\)
hacia \(+y\), en sentido antihorario. Entonces,
\[ v_x=v\cos\theta_v,\qquad v_y=v\sen\theta_v. \]
Si conocemos las componentes, \(\tan\theta_v=v_y/v_x\) cuando
\(v_x\ne0\), pero el cociente no distingue cuadrantes opuestos: hay que
utilizar los signos de ambas componentes. Si \(v_x=0\) y \(v_y\ne0\),
la velocidad es paralela al eje \(y\).

En tres dimensiones un único ángulo no basta. Para \(v>0\), los ángulos
\(\alpha\), \(\beta\) y \(\gamma\) respecto de los semiejes positivos
satisfacen
\[ \cos\alpha=\frac{v_x}{v},\qquad
\cos\beta=\frac{v_y}{v},\qquad \cos\gamma=\frac{v_z}{v}, \]
y \(\cos^2\alpha+\cos^2\beta+\cos^2\gamma=1\). No son tres datos
independientes. También podemos especificar la orientación mediante
el unitario \(\vect{v}/v\), sin calcular ángulos.

\begin{ejemplo}[Componente negativa y rapidez positiva]
En una mesa, con \(+x\) a la derecha y \(+y\) hacia el fondo, una partícula
tiene \(\vect{v}=(-\SI{3}{\meter\per\second})\ihat
+(\SI{4}{\meter\per\second})\jhat\). Su rapidez es
\(v=\SI{5}{\meter\per\second}\). Se mueve hacia valores menores de \(x\)
y mayores de \(y\); la velocidad apunta al segundo cuadrante:
\[ \theta_v=180\grado-\arctan(4/3)\approx126.9\grado. \]
El valor \(-53.1\grado\) que puede devolver una calculadora para
\(\arctan(-4/3)\) señalaría el cuarto cuadrante y sería incorrecto.
La dirección de esta velocidad no informa dónde está la partícula.
\end{ejemplo}

\section{De las gráficas temporales a la trayectoria}

Una gráfica \(x\)-\(t\) conserva la interpretación del capítulo 8: su
pendiente es \(v_x\). Las pendientes de \(y\)-\(t\) y \(z\)-\(t\)
proporcionan \(v_y\) y \(v_z\). Para reconstruir la posición hay que leer
todas las gráficas en el \emph{mismo instante}.

En cambio, una gráfica de la trayectoria \(y\)-\(x\) tiene dos ejes
espaciales. Su pendiente describe orientación, no rapidez. Si \(v_x\ne0\),
la regla de la cadena da
\[ \deriv{y}{x}=\frac{\dd y/\dd t}{\dd x/\dd t}=\frac{v_y}{v_x}. \]
La pendiente \(\dd y/\dd x\) es adimensional; \(\dd x/\dd t\) tiene
unidades de velocidad. Si \(v_x=0\) y \(v_y\ne0\), la tangente es
vertical en la gráfica \(y\)-\(x\) y el cociente no es finito.

\begin{ejemplo}[Una descripción en tres gráficas]
Un punto luminoso recorre una pantalla plana, con origen en su esquina
inferior izquierda, \(+x\) a la derecha y \(+y\) arriba. Durante
\(0\leq t\leq\SI{3}{\second}\), su programa establece
\[ x(t)=(\SI{2}{\meter\per\second})t,\qquad
y(t)=(\SI{0.5}{\meter\per\second\squared})t^2. \]
La coordenada \(x\) crece uniformemente; \(y\) aumenta cada vez más deprisa.
Al derivar,
\[ v_x=\SI{2}{\meter\per\second},\qquad
v_y=(\SI{1}{\meter\per\second\squared})t. \]
A \(t=\SI{2}{\second}\), la posición es
\((\SI{4}{\meter},\SI{2}{\meter})\) y
\[ \vect{v}=(\SI{2}{\meter\per\second})\ihat
+(\SI{2}{\meter\per\second})\jhat. \]
La rapidez vale \(2\sqrt{2}\,\si{\meter\per\second}\approx
\SI{2.83}{\meter\per\second}\), y la velocidad forma \(45\grado\)
con \(+x\). Eliminando el tiempo,
\[ y=\frac{x^2}{\SI{8}{\meter}},\qquad 0\leq x\leq\SI{6}{\meter}. \]
La pendiente de esta trayectoria en \(x=\SI{4}{\meter}\) es \(1\),
coherente con \(v_y/v_x=1\). El vector de posición apunta desde el origen
hasta \((4,2)\), pero la velocidad apunta a lo largo de la tangente:
no tienen igual orientación. La curva es parabólica por las funciones
programadas, sin que el punto luminoso sea un proyectil.
\end{ejemplo}

\begin{figure}[htbp]
    \centering
    \begin{tikzpicture}
        \begin{axis}[width=6.2cm,height=4.6cm,xmin=0,xmax=3.3,ymin=0,ymax=6.6,
            xlabel={$t\;(\si{\second})$},ylabel={$x\;(\si{\meter})$},
            grid=major,xtick={0,1,2,3},ytick={0,2,4,6}]
            \addplot[AzulPrincipal,thick,domain=0:3] {2*x};
            \addplot[only marks,AzulPrincipal] coordinates {(2,4)};
        \end{axis}
    \end{tikzpicture}
    \hfill
    \begin{tikzpicture}
        \begin{axis}[width=6.2cm,height=4.6cm,xmin=0,xmax=3.3,ymin=0,ymax=5,
            xlabel={$t\;(\si{\second})$},ylabel={$y\;(\si{\meter})$},
            grid=major,xtick={0,1,2,3},ytick={0,1,2,3,4}]
            \addplot[VerdePrincipal,thick,domain=0:3] {0.5*x^2};
            \addplot[only marks,VerdePrincipal] coordinates {(2,2)};
        \end{axis}
    \end{tikzpicture}

    \begin{tikzpicture}
        \begin{axis}[width=9cm,height=6.8cm,axis equal image,
            xmin=0,xmax=6.6,ymin=0,ymax=5,
            xlabel={$x\;(\si{\meter})$},ylabel={$y\;(\si{\meter})$},
            grid=major,xtick={0,1,2,3,4,5,6},ytick={0,1,2,3,4,5}]
            \addplot[AzulPrincipal,thick,domain=0:6,samples=60] {x^2/8};
            \addplot[only marks] coordinates {(0,0) (2,0.5) (4,2) (6,4.5)};
            \addplot[dashed,GrisPrincipal,domain=2.8:5.5] {x-2};
            \draw[-{Latex},very thick,VerdePrincipal]
                (axis cs:4,2) -- (axis cs:5,3) node[right] {$\vect{v}$};
            \node[anchor=north west] at (axis cs:4,1.9) {$t=\SI{2}{\second}$};
        \end{axis}
    \end{tikzpicture}
    \caption{Arriba, coordenadas en función del tiempo. Abajo, trayectoria
    con posiciones separadas por \(\SI{1}{\second}\) y tangente en
    \(t=\SI{2}{\second}\). La flecha indica la orientación de la velocidad;
    su longitud no se mide con la escala espacial.}
\end{figure}

\section{Aceleración: cambio de rapidez y cambio de dirección}

\subsection{La resta de velocidades es vectorial}

Un vehículo que toma una curva cambia su velocidad aunque su indicador de
rapidez mantenga la misma lectura. En una recta seguíamos el cambio de una
componente; ahora también puede cambiar la orientación del vector.
\[ \overline{\vect{a}}=\frac{\Delta\vect{v}}{\Delta t},
\qquad \Delta\vect{v}=\vect{v}_f-\vect{v}_i. \]
Por componentes,
\[ \bar a_x=\frac{v_{fx}-v_{ix}}{\Delta t},\quad
\bar a_y=\frac{v_{fy}-v_{iy}}{\Delta t},\quad
\bar a_z=\frac{v_{fz}-v_{iz}}{\Delta t}. \]
Para representar \(\Delta\vect{v}\), trasladamos las dos velocidades a
un origen común y dibujamos una flecha desde la punta de \(\vect{v}_i\)
hasta la de \(\vect{v}_f\). Es un diagrama de velocidades, no de posiciones.

\begin{ejemplo}[Aceleración media sin cambio neto de rapidez]
Un vehículo pasa de moverse al este con \(\SI{5}{\meter\per\second}\)
a moverse al norte con igual rapidez en \(\SI{2}{\second}\). Con \(+x\)
al este y \(+y\) al norte,
\[ \vect{v}_i=(\SI{5}{\meter\per\second})\ihat,
\qquad \vect{v}_f=(\SI{5}{\meter\per\second})\jhat. \]
Por tanto,
\[ \Delta\vect{v}=(-\SI{5}{\meter\per\second})\ihat
+(\SI{5}{\meter\per\second})\jhat, \]
\[ \overline{\vect{a}}=(-\SI{2.5}{\meter\per\second\squared})\ihat
+(\SI{2.5}{\meter\per\second\squared})\jhat. \]
Su magnitud es aproximadamente \(\SI{3.54}{\meter\per\second\squared}\),
y apunta al noroeste. Restar las rapideces habría dado cero y perdido la
información del giro. Los datos de los extremos no permiten conocer la
aceleración instantánea ni asegurar rapidez constante durante todo el intervalo.
\end{ejemplo}

\begin{figure}[htbp]
    \centering
    \begin{tikzpicture}[scale=0.7]
        \draw[-{Latex}] (0,0) -- (6.2,0) node[right] {$v_x\;(\si{\meter\per\second})$};
        \draw[-{Latex}] (0,0) -- (0,6.2) node[above] {$v_y\;(\si{\meter\per\second})$};
        \draw[-{Latex},very thick,AzulPrincipal] (0,0) -- (5,0)
            node[midway,below] {$\vect{v}_i$};
        \draw[-{Latex},very thick,AzulPrincipal] (0,0) -- (0,5)
            node[midway,left] {$\vect{v}_f$};
        \draw[-{Latex},very thick,RojoPrincipal] (5,0) -- (0,5)
            node[midway,above right] {$\Delta\vect{v}$};
        \node[below] at (5,0) {$5$};
        \node[left] at (0,5) {$5$};
        \node[below left] at (0,0) {$0$};
    \end{tikzpicture}
    \caption{Resta de velocidades del ejemplo. Las flechas azules tienen
    igual magnitud, pero su diferencia no es nula.}
\end{figure}

\subsection{Aceleración instantánea por componentes}

Aplicamos el límite del capítulo 9 al vector completo:
\[ \vect{a}=\lim_{\Delta t\to0}\frac{\Delta\vect{v}}{\Delta t}
=\deriv{\vect{v}}{t}=\derivdos{\vect{r}}{t}. \]
En nuestros ejes cartesianos fijos,
\[
\begin{aligned}
\vect{a}&=a_x\ihat+a_y\jhat+a_z\khat,\\
a_x&=\deriv{v_x}{t}=\derivdos{x}{t},\qquad
a_y=\deriv{v_y}{t}=\derivdos{y}{t},\qquad
a_z=\deriv{v_z}{t}=\derivdos{z}{t}.
\end{aligned}
\]
Su magnitud es \(a=\sqrt{a_x^2+a_y^2+a_z^2}\), y la unidad SI es
\(\si{\meter\per\second\squared}\). La aceleración no tiene que ser
paralela a la velocidad ni apuntar al origen. Para el punto luminoso
anterior, \(a_x=0\) y \(a_y=\SI{1}{\meter\per\second\squared}\):
la aceleración es vertical mientras la velocidad cambia de orientación.

\subsection{¿Qué parte de la aceleración modifica la rapidez?}

Para \(v>0\), introducimos el unitario tangente \(\uvec{t}=\vect{v}/v\).
Al escribir \(\vect{v}=v\uvec{t}\) y derivar un producto,
\[ \vect{a}=\deriv{v}{t}\uvec{t}+v\deriv{\uvec{t}}{t}. \]
El primer término cambia la rapidez; el segundo cambia la orientación.
Como \(\uvec{t}\pdot\uvec{t}=1\), su derivada implica
\(\uvec{t}\pdot\dd\uvec{t}/\dd t=0\): el cambio del unitario es
perpendicular a él. La componente de aceleración paralela a la velocidad
es, por tanto, la que determina la variación de rapidez.

El producto escalar permite medir esa proyección. Para dos vectores
\(\vect{A}\) y \(\vect{B}\), se define en esta base ortonormal como
\(\vect{A}\pdot\vect{B}=A_xB_x+A_yB_y+A_zB_z\); si ambos son no nulos,
también equivale a \(AB\cos\phi\), donde \(\phi\) es el ángulo entre
ellos. Es una cantidad escalar, positiva para ángulos agudos, negativa
para obtusos y nula para vectores perpendiculares.

Podemos expresar el cambio de rapidez mediante este producto:
\[ v^2=\vect{v}\pdot\vect{v}
\quad\Longrightarrow\quad 2v\deriv{v}{t}=2\vect{v}\pdot\vect{a}. \]
Así, para \(v\ne0\),
\[ \deriv{v}{t}=\frac{\vect{v}\pdot\vect{a}}{v}
=\frac{v_xa_x+v_ya_y+v_za_z}{v}. \]
Si el producto escalar es positivo, la rapidez aumenta; si es negativo,
disminuye. Si es cero, su tasa instantánea de cambio es cero, aunque la
dirección puede estar cambiando. Que se anule en un instante no implica
rapidez constante durante un intervalo.

Esta relación generaliza el criterio de signos del movimiento rectilíneo.
No basta examinar \(a_x\) o \(a_y\) por separado para saber si aumenta
la rapidez total. Tampoco se identifica \(a\) con \(\dd v/\dd t\): la
magnitud de la aceleración incluye el efecto del cambio de dirección.
En reposo instantáneo no se divide por \(v\); se examinan directamente
las componentes y su evolución próxima. La relación entre aceleración
normal y radio de curvatura se retomará al estudiar movimiento circular.

\section{Componentes de un mismo movimiento}

\subsection{Qué significa analizar los ejes por separado}

Una igualdad vectorial en una base cartesiana equivale a tres igualdades
escalares. De este modo,
\[ \deriv{\vect{v}}{t}=\vect{a}
\quad\Longleftrightarrow\quad
\deriv{v_x}{t}=a_x,\quad \deriv{v_y}{t}=a_y,\quad \deriv{v_z}{t}=a_z. \]
La componente \(a_y\) no se suma a \(a_x\) para calcular el cambio de
\(v_x\). Si \(a_x=0\) durante un intervalo, \(v_x\) permanece constante
en él, aunque cambien \(v_y\) y \(v_z\). Esto no exige que \(v_x=0\).

Las proyecciones no son etapas sucesivas del recorrido: la partícula
evoluciona en todos los ejes simultáneamente. Debemos usar el mismo reloj
y las condiciones iniciales del mismo instante. No existen tiempos de
movimiento horizontal y vertical independientes para un único evento.

\begin{advertencia}[Descomposición e independencia]
La descomposición cartesiana es válida bajo la elección de ejes fijos
utilizada aquí. Resolver las ecuaciones \emph{independientemente} requiere
además que los datos o el modelo lo permitan. Si \(a_x\) depende de
\(y(t)\), las ecuaciones están acopladas: para conocer \(x(t)\) puede
ser necesario conocer también \(y(t)\). Con aceleración constante, o con
cada componente de aceleración prescrita solo como función del tiempo,
podemos integrar los ejes por separado.
\end{advertencia}

\subsection{Reconstrucción a partir de velocidad o aceleración}

Derivar permite pasar de posición a velocidad y de velocidad a aceleración.
El proceso inverso acumula pequeños cambios a lo largo del tiempo.
Si conocemos \(\vect{v}(t)\),
\[ \vect{r}(t)=\vect{r}_0+\int_{t_0}^{t}\vect{v}(t')\,\dd t'. \]
El símbolo \(t'\) es una variable de integración. Por ejemplo,
\[ x(t)=x_0+\int_{t_0}^{t}v_x(t')\,\dd t', \]
y se procede igual para \(y\) y \(z\). El área algebraica bajo
\(v_x\)-\(t\) es \(\Delta x\), no la distancia total. Las tres áreas
proporcionan las componentes de \(\Delta\vect{r}\).

Si conocemos la aceleración y la velocidad inicial,
\[ \vect{v}(t)=\vect{v}_0+\int_{t_0}^{t}\vect{a}(t')\,\dd t'. \]
Cada área algebraica bajo \(a_x\)-\(t\), \(a_y\)-\(t\) o \(a_z\)-\(t\)
es el cambio de la componente correspondiente de velocidad. Hace falta
\(\vect{v}_0\) para determinar la velocidad y, además, \(\vect{r}_0\)
para determinar la posición: la aceleración sola no especifica el movimiento.

Para la distancia recorrida desde \(t_0\), con \(s(t_0)=0\), acumulamos
la rapidez:
\[ s(t)=\int_{t_0}^{t}v(t')\,\dd t'
=\int_{t_0}^{t}\sqrt{v_x(t')^2+v_y(t')^2+v_z(t')^2}\,\dd t'. \]
Primero se calcula la magnitud y luego se integra. Integrar el vector y
tomar después su norma daría la magnitud del desplazamiento, en general
menor. Las unidades lo confirman: velocidad por tiempo es longitud,
y aceleración por tiempo es velocidad.

\section{Movimiento con aceleración vectorial constante}

\subsection{Hipótesis y construcción de las ecuaciones}

Supongamos que \(\vect{a}\) permanece constante en el intervalo estudiado.
Sus componentes sobre nuestros ejes fijos deben ser constantes; conservar
solo la magnitud no basta. Definimos \(\Delta t=t-t_0\). Al integrar,
\[ \vect{v}(t)=\vect{v}_0+\vect{a}\Delta t. \]
La velocidad es la suma de la velocidad inicial y el cambio acumulado
\(\vect{a}\Delta t\). Integrando de nuevo,
\[
\begin{aligned}
\vect{r}(t)-\vect{r}_0
&=\int_{t_0}^{t}\bigl[\vect{v}_0+\vect{a}(t'-t_0)\bigr]\,\dd t'\\
&=\vect{v}_0\Delta t+\frac12\vect{a}(\Delta t)^2.
\end{aligned}
\]
El primer término es el desplazamiento que habría con la velocidad inicial
constante; el segundo es la corrección debida a la aceleración. Son
contribuciones simultáneas, no dos tramos recorridos por separado.

Las ecuaciones equivalentes por componentes son
\[
\begin{aligned}
v_x&=v_{0x}+a_x\Delta t,&
x&=x_0+v_{0x}\Delta t+\tfrac12a_x(\Delta t)^2,\\
v_y&=v_{0y}+a_y\Delta t,&
y&=y_0+v_{0y}\Delta t+\tfrac12a_y(\Delta t)^2,\\
v_z&=v_{0z}+a_z\Delta t,&
z&=z_0+v_{0z}\Delta t+\tfrac12a_z(\Delta t)^2.
\end{aligned}
\]
Si \(t_0=0\), \(\Delta t=t\), y recuperamos en cada eje las expresiones
del capítulo 9. En cada ecuación de posición todos los términos tienen
dimensiones de longitud. Si la aceleración es aproximadamente constante,
estas expresiones también son aproximadas y se restringen al intervalo
en que el modelo resulta adecuado.

\subsection{Consecuencias y casos particulares}

Si \(\vect{a}=\vect{0}\) durante todo el intervalo, la velocidad es
constante. Para \(\vect{v}_0\ne\vect{0}\) la trayectoria es recta aunque
varíen varias coordenadas: es el MRU expresado con vectores. Si también
\(\vect{v}_0=\vect{0}\), la partícula permanece en reposo.

Con aceleración constante no nula, si \(\vect{v}_0\) es paralela a
\(\vect{a}\), o nula, el movimiento sigue una recta. Si no son paralelos,
todos los desplazamientos son combinaciones de esos dos vectores: la
trayectoria queda en el plano que ambos generan, trasladado hasta
\(\vect{r}_0\). Usar tres componentes no impide que el movimiento sea
plano. La trayectoria de los proyectiles se estudiará en el próximo capítulo.

La linealidad de \(\vect{v}(t)\) permite escribir, bajo esta misma hipótesis,
\[ \overline{\vect{v}}=\frac{\vect{v}_0+\vect{v}(t)}{2},\qquad
\Delta\vect{r}=\frac{\vect{v}_0+\vect{v}(t)}{2}\,\Delta t. \]
No se deduce que la rapidez media sea la semisuma de las rapideces:
tomar la magnitud no es una operación lineal.

La relación sin tiempo también se generaliza con cuidado. Multiplicando
escalarmente el desplazamiento por \(2\vect{a}\),
\[ 2\vect{a}\pdot\Delta\vect{r}
=2\vect{a}\pdot\vect{v}_0\Delta t+a^2(\Delta t)^2. \]
El miembro derecho coincide con la diferencia entre
\(\norm{\vect{v}_0+\vect{a}\Delta t}^2\) y \(v_0^2\). Por tanto,
\[ v^2=v_0^2+2\vect{a}\pdot\Delta\vect{r}
=v_0^2+2(a_x\Delta x+a_y\Delta y+a_z\Delta z). \]
Aquí \(v\) y \(v_0\) son rapideces. No se sustituye el producto escalar
por \(as\): la distancia recorrida no conserva la orientación.
Esta ecuación escalar tampoco determina por sí sola las tres componentes
finales de velocidad.

\begin{ejemplo}[Una proyección se invierte mientras la partícula avanza]
En una mesa horizontal un punto controlado por un mecanismo parte en
\(t_0=0\) de \(x_0=\SI{1}{\meter}\), \(y_0=\SI{2}{\meter}\), con
\(v_{0x}=\SI{2}{\meter\per\second}\) y
\(v_{0y}=-\SI{2}{\meter\per\second}\). Tomamos \(+x\) a la derecha,
\(+y\) hacia el fondo. Durante \(0\leq t\leq\SI{4}{\second}\),
\(a_x=0\) y \(a_y=\SI{1}{\meter\per\second\squared}\).

La componente \(x\) corresponde a un MRU; la componente \(y\), a un MRUA:
\[
\begin{aligned}
x(t)&=\SI{1}{\meter}+(\SI{2}{\meter\per\second})t,\\
y(t)&=\SI{2}{\meter}-(\SI{2}{\meter\per\second})t
+\tfrac12(\SI{1}{\meter\per\second\squared})t^2,\\
v_x(t)&=\SI{2}{\meter\per\second},\\
v_y(t)&=-\SI{2}{\meter\per\second}
+(\SI{1}{\meter\per\second\squared})t.
\end{aligned}
\]
En \(t=\SI{2}{\second}\), \(y\) alcanza su mínimo y \(v_y=0\), pero
\(v_x\ne0\). La partícula pasa por \((\SI{5}{\meter},0)\) con rapidez
\(\SI{2}{\meter\per\second}\) y tangente paralela a \(x\); no se detiene.
Antes, \(v_ya_y<0\) y la rapidez disminuye. Después, \(v_ya_y>0\) y
aumenta, aunque la aceleración no haya cambiado.

En \(t=\SI{4}{\second}\), la posición es
\((\SI{9}{\meter},\SI{2}{\meter})\), y
\(\vect{v}=(\SI{2}{\meter\per\second})\ihat
+(\SI{2}{\meter\per\second})\jhat\). Regresa al valor inicial de \(y\),
pero no a la posición inicial: \(\Delta\vect{r}=(\SI{8}{\meter})\ihat\).
La velocidad media es \((\SI{2}{\meter\per\second})\ihat\), en acuerdo
con la semisuma vectorial de las velocidades extremas. Además,
\(\vect{a}\pdot\Delta\vect{r}=0\): las rapideces inicial y final son
iguales, aunque la rapidez intermedia haya sido menor.
\end{ejemplo}

\section{Una trayectoria genuinamente espacial}

La descripción en tres dimensiones no exige nuevas definiciones, pero sí
conservar las tres componentes al reconstruir vectores y magnitudes.
Un movimiento con aceleración variable puede abandonar cualquier plano fijo.

\begin{ejemplo}[Posición, velocidad y aceleración en el espacio]
Un mecanismo desplaza una sonda respecto de un origen fijo del laboratorio.
Elegimos \(+x\) al este, \(+y\) al norte y \(+z\) hacia arriba.
Durante \(0\leq t\leq\SI{2}{\second}\), su movimiento se modela mediante
\[
\begin{aligned}
x(t)&=(\SI{1}{\meter\per\second})t,\\
y(t)&=(\SI{0.5}{\meter\per\second\squared})t^2,\\
z(t)&=\SI{1}{\meter}+(\SI{0.25}{\meter\per\second\cubed})t^3.
\end{aligned}
\]
Los coeficientes tienen las unidades necesarias para obtener longitudes.
La sonda avanza hacia el este, el norte y arriba; no se supone caída libre.
Derivando cada coordenada,
\[
\begin{aligned}
\vect{v}(t)&=(\SI{1}{\meter\per\second})\ihat
+(\SI{1}{\meter\per\second\squared})t\jhat
+(\SI{0.75}{\meter\per\second\cubed})t^2\khat,\\
\vect{a}(t)&=(\SI{1}{\meter\per\second\squared})\jhat
+(\SI{1.5}{\meter\per\second\cubed})t\khat.
\end{aligned}
\]
La aceleración no es constante; no sería válido insertar su valor final
en las ecuaciones del apartado anterior para todo el intervalo.

En \(t=\SI{2}{\second}\),
\[
\begin{aligned}
\vect{r}&=(2\ihat+2\jhat+3\khat)\,\si{\meter},\\
\vect{v}&=(\ihat+2\jhat+3\khat)\,\si{\meter\per\second},\\
\vect{a}&=(\jhat+3\khat)\,\si{\meter\per\second\squared}.
\end{aligned}
\]
La rapidez es \(\sqrt{14}\,\si{\meter\per\second}\approx
\SI{3.74}{\meter\per\second}\); la magnitud de la aceleración es
\(\sqrt{10}\,\si{\meter\per\second\squared}\approx
\SI{3.16}{\meter\per\second\squared}\). La dirección de la velocidad
queda especificada por \((\ihat+2\jhat+3\khat)/\sqrt{14}\).
Además,
\[ \vect{v}\pdot\vect{a}=\SI{11}{\meter\squared\per\second\cubed}>0, \]
por lo que la rapidez aumenta. Su tasa de cambio es
\(11/\sqrt{14}\,\si{\meter\per\second\squared}\approx
\SI{2.94}{\meter\per\second\squared}\), distinta de \(a\).

Al comienzo, \(\vect{r}_0=(\SI{1}{\meter})\khat\). En dos segundos,
\(\Delta\vect{r}=(2\ihat+2\jhat+2\khat)\,\si{\meter}\), de modo que
\(\overline{\vect{v}}=(\ihat+\jhat+\khat)\,\si{\meter\per\second}\).
La posición inicial no puede omitirse al calcular el desplazamiento.

Esta curva no está contenida en un plano. Si existieran constantes
\(A,B,C,D\), con \(A,B,C\) no todas nulas, tales que \(Ax+By+Cz=D\)
durante todo el intervalo, al sustituir obtendríamos una identidad con
términos en \(t\), \(t^2\) y \(t^3\). Sus coeficientes deben anularse
por separado: resultaría \(A=B=C=0\), contradiciendo la condición de
un plano. Aquí la tercera dimensión es esencial.
\end{ejemplo}

\section{Cómo observar y comprobar la descripción}

Registrar posiciones a intervalos iguales permite relacionar los vectores
con una observación real. Para un movimiento plano puede filmarse un punto
desde una cámara fija, con el eje óptico perpendicular al plano del
movimiento y una escala de longitud en ese mismo plano. Después se asignan
coordenadas a cada imagen y tiempos mediante el registro de la cámara.
La perspectiva y la calibración deben controlarse antes de interpretar
píxeles como distancias.

Con un intervalo de muestreo \(h\), podemos estimar, en instantes interiores,
\[ \vect{v}(t)\approx\frac{\vect{r}(t+h)-\vect{r}(t-h)}{2h},\qquad
\vect{a}(t)\approx\frac{\vect{r}(t+h)-2\vect{r}(t)+\vect{r}(t-h)}{h^2}. \]
Son aproximaciones para movimientos suficientemente suaves, no definiciones
exactas con \(h\) finito. Reducir \(h\) mejora la resolución temporal,
pero puede hacer más importante el error de localización, especialmente
en la segunda diferencia usada para la aceleración. Una sola imagen plana
no determina en general la profundidad de una trayectoria espacial: hacen
falta datos adicionales, como vistas calibradas desde otras posiciones.

\begin{estrategia}[Leer el movimiento antes de calcular]
Fije el origen, los ejes y el reloj; identifique si conoce posición,
velocidad o aceleración. Trabaje por componentes con sus signos y unidades,
manteniendo el mismo instante para todas ellas. Use derivadas para obtener
tasas de cambio e integrales con condiciones iniciales para reconstruir el
movimiento. Aplique las ecuaciones de aceleración constante solo después
de comprobar esa hipótesis. Al final, reconstruya el vector, su magnitud
y su orientación, y contraste el resultado con la trayectoria y las gráficas.
\end{estrategia}

\section{Comprobación conceptual y continuidad}

\begin{errorfrecuente}[Confundir la proyección con el movimiento completo]
Una componente nula no hace nulo al vector. Por ejemplo, \(v_y=0\) no
implica reposo si \(v_x\ne0\). Del mismo modo, \(a_x=0\) permite que
\(v_x\) sea una constante no nula. Volver a la coordenada inicial
\(y_0\) tampoco equivale a volver al punto inicial: deben coincidir
todas las coordenadas en el mismo instante.
\end{errorfrecuente}

Como comprobación de comprensión, conviene discutir estas preguntas:
\begin{itemize}
    \item Si dos partículas recorren la misma curva, ¿qué falta para
    afirmar que tienen el mismo movimiento? La dependencia temporal
    de sus posiciones.
    \item Si la rapidez es constante, ¿debe ser nula la aceleración?
    No: la orientación de la velocidad puede variar.
    \item Si \(\vect{a}\) es constante y tiene tres componentes no nulas,
    ¿debe la trayectoria abandonar todo plano? No: el desplazamiento
    pertenece al plano generado por \(\vect{v}_0\) y \(\vect{a}\),
    o a una recta en el caso degenerado.
\end{itemize}

\begin{notapedagogica}
Para presentar el tema puede seguirse un punto sobre una mesa: registrar
posiciones, unirlas mediante desplazamientos y comparar después las
velocidades tangentes en puntos sucesivos. Contrastar la trayectoria con
las gráficas temporales separa geometría y evolución temporal. El ejemplo
de la proyección que invierte su sentido muestra por qué los ejes deben
analizarse por separado y reconstruirse conjuntamente.
\end{notapedagogica}

\begin{resumen}
Posición, velocidad y aceleración son vectores distintos: describen dónde
está la partícula, cómo cambia su posición y cómo cambia su velocidad.
Las componentes cartesianas permiten aplicar las relaciones conocidas de
una dimensión a cada proyección, siempre con un tiempo común. La velocidad
no nula es tangente a la trayectoria; la aceleración puede modificar su
magnitud, su orientación o ambas. Las magnitudes se reconstruyen a partir
de todas las componentes; la distancia se obtiene acumulando la rapidez.
\end{resumen}

El movimiento vertical anterior se recupera al mantener fijas las
coordenadas horizontales. En el próximo capítulo combinaremos esa
descripción vertical con el movimiento horizontal para estudiar proyectiles
bajo hipótesis explícitas. El movimiento circular permitirá cuantificar
el cambio de dirección, y el movimiento relativo comparará descripciones
realizadas por observadores diferentes.
